\documentclass{article}
\usepackage{tcee-repaso}
\tcetema[Teoría de la integración económica]{3.B.10}{Teoría de la integración económica.}
\tcefecha{05/11/2026}
\begin{document}
\cabecerarepaso
\begin{columnas}

% ==========================================================================
\bloque{Esquema del cante}
\esquemacante{%
\apartadocante{Introducción: «dime con quién andas\ldots»}{2,5}
\apartadocante{I. Efectos de la integración económica}{13}
\subcante{1. Formas de integración y enfoques}{1}
\subcante{2. Efectos estáticos de la UA}{6}
\subcante{3. ¿Por qué puede ser mejor la UA?}{5}
\subcante{4. Formas superiores}{1}
\apartadocante{II. Determinantes de la integración}{5}
\subcante{Advertencia: las 2 caras de la moneda}{0,25}
\subcante{1. Determinantes geográficos}{2,25}
\subcante{2. Determinantes políticos}{2,5}
\apartadocante{III. Metodología empírica y evidencia}{8}
\subcante{Idea: \emph{ex ante} vs. \emph{ex post}}{0,25}
\subcante{1. Modelos de equilibrio general computable}{1,25}
\subcante{2. La ecuación de gravedad}{3,75}
\subcante{3. Nuevos modelos cuantitativos}{0,75}
\subcante{4. Evidencia}{2}
\apartadocante{Conclusión}{1,5}}
\emph{Pizarra} en 3 columnas: \hyperref[rep:I]{I} (escalera y \autor{Viner}-\autor{Lipsey} por capas), \hyperref[rep:II]{II} (\emph{tomahawk}), \hyperref[rep:III]{III} (gravedad). Pegatinas: «para este viaje no se necesitaban alforjas» (\autor{Cooper}-\autor{Massell}), «quien da primero da dos veces» (especialización perversa), «arrimarse al sol que más calienta» (centro-periferia), «cuando las barbas de tu vecino veas pelar\ldots» (dominó), «el roce hace el cariño» (gravedad).

% ==========================================================================
\bloque[2,5]{Introducción}
\begin{itemize}
\item \textbf{Enganche}: «Dime con quién andas y te diré quién eres»: el efecto de una UA depende de \emph{con quién} se forma. UE-Mercosur, aplicación provisional desde el 1/5/2026: ¿proveedor más barato o desviación de comercio? ¿Mejora el bienestar suprimir barreras con unos y mantenerlas frente a otros?
\item \textbf{Relevancia}: \emph{práctica}, 390 ACR notificados a la OMC y en vigor ($\approx$9 de cada 10, ZLC) + $\geq$84 sin notificar; solo 16 UA: nadie quiere ceder la política comercial. \emph{Teórica}: el caso más claro de \emph{segundo óptimo}.
\item \textbf{Contextualización}: \concepto{integración económica} = eliminación de las «fronteras económicas» (todo lo que restrinja la circulación de bienes, servicios y factores) entre $\geq$2 países (\autor{Pelkmans}); cara \emph{negativa} (suprimir obstáculos) y \emph{positiva} (instituciones y políticas comunes) (\autor{Tinbergen}, 1954). \autor{Balassa} (1961): \emph{proceso} y \emph{estado}. Siempre \emph{discriminatoria} ($\neq$ liberalización unilateral o multilateral). Historia: tradicional (50-60, con la CEE); nueva teoría (80: escala, competencia imperfecta, crecimiento); NGE (1991); economía política y gravedad (90). Segundo óptimo del librecambio de \autor{Smith} y \autor{Ricardo} \vertema{3.A.22}.
\item \textbf{Problemática}: ¿mejora el bienestar si discrimina? ¿Por qué y con quién se integran los países? ¿Cómo se mide y qué dicen los datos? ¿Peldaño o escollo?
\item \textbf{Estructura}: I. Efectos (estático clásico y lo que deja fuera); II. Determinantes (geografía y política); III. Medición y evidencia (gravedad).
\end{itemize}

% ==========================================================================
\bloque[13]{I. Efectos de la integración económica\label{rep:I}}
\begin{ideaclave}
Efectos estáticos de una UA ambiguos y dominados por una rebaja no discriminatoria; la justifican compensaciones, RRI, escala, competencia y crecimiento, con costes redistributivos.
\end{ideaclave}

\sub{I.1. Formas de integración y enfoques\label{rep:I.1}}
\begin{itemize}
\item \textbf{Escalera} (\autor{Balassa}, 1961, + mercado único): preferencias (rebaja parcial) $\to$ \concepto{ZLC} (aranceles 0 entre socios) $\to$ \concepto{UA} (+ arancel exterior común) $\to$ \concepto{mercado común} (+ factores) $\to$ \concepto{mercado único} (+ barreras no arancelarias) $\to$ \concepto{unión económica} (+ coordinación de políticas; y moneda en la UEM) $\to$ \concepto{integración total} (+ instituciones supranacionales).
\item Ejemplos: D-8; T-MEC, AELC; CEE (UA 1968), SACU (1910); Mercosur (imperfecto); UE (1993); euro (21).
\item \concepto{Desviación del tráfico} (\emph{trade deflection}): en la ZLC se entra por el socio con arancel más bajo; se combate con \emph{reglas de origen} (encarecen: «plato de espaguetis»). La UA lo evita con arancel común a cambio de ceder la política comercial. $\neq$ desviación de comercio de \autor{Viner}.
\end{itemize}
\graficorepaso[0.8\linewidth]{zlc-ua}{ZLC con desviación del tráfico (izq.) vs. UA con arancel exterior común (dcha.)}
\begin{itemize}
\item \textbf{OMC}: choca con la NMF (art.~I GATT) \vertema{3.B.33}. Art.~XXIV: ZLC y UA que cubran «lo esencial de los intercambios»; en la UA, arancel común no más alto «en conjunto» (promedio ponderado desde el Entendimiento de 1994). Preferencias parciales solo entre PED (cláusula de habilitación, 1979); servicios, art.~V AGCS \vertema{3.B.34}.
\item Pocas formas puras; acuerdos de \emph{nueva generación} (servicios, inversión, contratación pública, normas). 44 ACR en vigor hasta 1990; +145 en los 90; salto de 2021 por el \emph{Brexit} (38 de 46).
\end{itemize}
\graficorepaso[0.85\linewidth]{acr-en-vigor}{ACR notificados al GATT/OMC y en vigor por tipo, 1958-2026 (OMC, RTA-IS)}
\begin{itemize}
\item \textbf{2 enfoques}: \emph{tradicional} (estático, competencia perfecta) y \emph{nueva teoría} (escala, competencia imperfecta, crecimiento, localización).
\end{itemize}

\sub{I.2. Efectos estáticos: país pequeño, competencia perfecta\label{rep:I.2}}
\begin{itemize}
\item \textbf{\autor{Viner} (1950)}: UA = proteccionismo \emph{erga extra} + libre comercio \emph{erga intra}. Se creía que siempre mejoraba; primero en rebatirlo (antes del segundo óptimo).
  \begin{itemize}
  \item \concepto{Creación de comercio}: se deja de producir caro en casa y se importa del socio, más eficiente.
  \item \concepto{Desviación de comercio}: se deja de importar del tercero barato para importar del socio caro, solo porque el arancel ya no le grava.
  \end{itemize}
\item \emph{Supuestos}: equilibrio parcial, $A$, $B$, $R$, un bien; $A$ pequeño, costes constantes $p_A$, $p_B$, $p_R$; arancel \emph{ad valorem} $t$; demanda \emph{perfectamente inelástica} $\bar x$; competencia perfecta.
\item \emph{Ejemplo} con $t=50$\,\%: (1) $B$ caro (16; $A$ 10; $R$ 12) $\to$ nada; (2) $B$ 8, $A$ 10, $R$ 12 $\to$ creación, gana $(p_A-p_B)\bar x=2\bar x$; (3) $B$ 8, $A$ 10, $R$ 6 $\to$ desviación:
  \begin{formula}
  \[\Delta W=[p_R(1+t)-p_B]\,\bar x-t\,p_R\,\bar x=(p_R-p_B)\,\bar x<0\]
  \end{formula}
  pierde lo que el socio es más caro que el eficiente.
\item Efecto \emph{ambiguo}: depende de con quién. $A$ pequeño: lo que gana $A$ lo gana el mundo. \emph{Crítica}: demanda fija.
\end{itemize}
\graficorepaso[\linewidth]{viner-situaciones}{\autor{Viner} con demanda inelástica: sin efecto, creación y desviación}
\begin{itemize}
\item \textbf{\autor{Meade} (1955)}: (i) incorpora la demanda: \concepto{expansión del comercio}; (ii) anticipa el \concepto{segundo óptimo} (\autor{Lipsey}-\autor{Lancaster}, 1956) \vertema{3.A.22}: si una condición de \autor{Pareto} no se cumple, no es necesariamente deseable cumplir las demás. UA $<$ libre comercio (con arancel a terceros se rompe $\mathrm{RMS}=\mathrm{RMT}=\mathrm{RRI}$) y no necesariamente $>$ arancel no discriminatorio: cuestión empírica.
\item \textbf{\autor{Lipsey} (1957)}, equilibrio parcial: demanda con \emph{pendiente negativa} (\concepto{creación de consumo}) y CMg \emph{creciente}; recaudación devuelta. $p_R<p_B<p_R(1+t)<p_B(1+t)$: con la UA el precio baja a $p_B$, $q_0\to q_1$, $x_0\to x_1$, todo de $B$. EC $+(a+b+c+d)$; EP $-a$; recaudación $-(c+e)$:
  \funcionclave{\Delta W=\underbrace{b}_{\text{creación de comercio}}+\underbrace{d}_{\text{creación de consumo}}-\underbrace{e}_{\text{desviación}}}
  Gana más con oferta y demanda elásticas ($b$, $d$ ↑) y socio poco más caro ($e$ ↓).
\end{itemize}
\graficorepaso[0.78\linewidth]{viner-lipsey}{UA (áreas $a$-$e$) y rebaja no discriminatoria hasta $p_R(1+t')=p_B$ (áreas $f$ y $g$)}
\begin{itemize}
\item \emph{Equilibrio general}: \emph{la desviación puede mejorar}: el arancel no discriminatorio distorsiona también el consumo ($\mathrm{RMS}$ = precio interior $\neq$ RRI); la UA suprime esa distorsión frente al socio $\Rightarrow$ curva de indiferencia más alta posible.
\item \textbf{Regla de \autor{Lipsey} (1960)}: más probable que mejore cuanto ↑ comercio con el socio, ↓ con el resto del mundo y ↓ comercio exterior / consumo nacional.
\end{itemize}
\begin{itemize}
\item \textbf{Factores del saldo} (más probable que mejore): \emph{antes}: aranceles previos altos; diferencia de costes $A$-$B$ grande y $B$-$R$ pequeña; poco comercio previo con $R$. \emph{Después}: arancel común bajo; oferta y demanda elásticas. \emph{Estructura}: más miembros; economías parecidas y productos sustitutivos (compiten: más creación); nº de bloques \emph{no monótono}: mínimo del bienestar mundial con $\approx$3 (\autor{Krugman}, 1991).
\item \concepto{Socios comerciales naturales} (\autor{Wonnacott}-\autor{Lutz}, 1989; \autor{Krugman}, \autor{Summers}, 1991): vecinos que ya comercian crearían sin desviar. \autor{Bhagwati}-\autor{Panagariya} (1996): no; si el socio es caro, más comercio previo = más recaudación perdida. \autor{Krishna} (2003): ni cercanía ni volumen predicen ganancias.
\item \textbf{\autor{Cooper}-\autor{Massell} (1965)}: rebaja no discriminatoria hasta $t'$ con $p_R(1+t')=p_B$: mismo precio, importa de $R$. EC y EP como la UA; recaudación de $c+e$ a $e+f+g$. Efecto neto $b+d+f+g>0$, inequívoco; gana $e+f+g$ más que la UA. ¿Por qué UA? Por los supuestos (pequeño, un mercado, competencia perfecta, estático).
\item \autor{Wonnacott}-\autor{Wonnacott} (1981): si el socio tiene aranceles y hay costes de transporte con $R$, la UA da \emph{acceso preferente al mercado del socio}: intercambio de concesiones. \autor{Cooper}-\autor{Massell} y \autor{Johnson} (1965): la producción industrial como bien público.
\end{itemize}

\sub{I.3. ¿Por qué puede ser mejor la UA?\label{rep:I.3}}
\begin{itemize}
\item \textbf{\autor{Kemp}-\autor{Wan} (1976)} (a partir de \autor{Ohyama}, 1972): \emph{siempre existe} una UA que mejora a los miembros sin empeorar a $R$: arancel común que deja \emph{inalterado el comercio con terceros} + transferencias de suma fija dentro. Sin desviación $\Rightarrow$ mejora de \autor{Pareto}. Límites: resultado de existencia; compensaciones rara vez pagadas (\autor{Bhagwati}, 1991). Ampliar uniones así lleva, paso a paso, al libre comercio mundial. $\neq$ art.~XXIV.5 (mira el \emph{nivel}; el teorema, el \emph{comercio} con terceros).
\item \textbf{País grande (\autor{Mundell}, 1964)}: RRI entre $A$ y $B$ ambigua; RRI de la unión frente a $R$ mejora \emph{inequívocamente} (importa menos, el precio baja): \emph{arancel óptimo} \vertema{3.B.7}, a costa de $R$ y con riesgo de represalias. Contrario a \autor{Kemp}-\autor{Wan}: mejora \emph{porque} perjudica. Más peso negociador \vertema{3.B.5}.
\item \textbf{Efectos procompetitivos}: ↓ poder de mercado. \autor{Brander}-\autor{Krugman} (1983): 2 monopolios pasan a \autor{Cournot} en los 2 mercados \vertema{3.A.19}: $q$ ↑, $p$ ↓; con costes de transporte altos el «dumping recíproco» puede restar \vertema{3.B.6}. Tipos: efecto precio (márgenes), no precio (innovación), ↓ ineficiencia X (\autor{Leibenstein}, 1966) \vertema{3.A.12}.
\item \textbf{Economías de escala, \autor{Corden} (1972)}: siguen creación y desviación y se añaden \concepto{reducción de costes} (+: lo que ya se producía, más barato) y \concepto{supresión de comercio} ($-$: nacen empresas que solo viven de la protección común). Supuestos: un productor como mucho por país, $p=\mathrm{CMe}$ (la CMe hace de oferta), demandas idénticas. Caso a), $p_R<p_B<p_A<p_R(1+t)$: $B$ abastece a los 2, $p^{\mathrm{UA}}$ donde $\mathrm{CMe}_B$ corta $D_{A+B}$; creación en $A$ y reducción de costes en $B$.
\end{itemize}
\pareja{corden}{\autor{Corden}, caso a): $B$ abastece la unión a $p^{\mathrm{UA}}$}{especializacion-perversa}{Especialización perversa: se queda la empresa menos eficiente}
\begin{itemize}
\item \emph{Casos}: a) 2$\to$1 productores: creación + reducción de costes, \emph{positivo}; b) 1$\to$1 (en $B$): desviación + reducción, \emph{ambiguo}; c) 0$\to$1: desviación + supresión, \emph{negativo en la producción}. Regla: menos productores ↑ bienestar; los mismos, ambiguo; más, se encarece.
\item \concepto{Especialización perversa}: sin demandas idénticas, la empresa con CMe más alta pero que parte de un coste más bajo (mercado mayor, empezó antes: \emph{accidente histórico}) se queda con el mercado ampliado. Misma dependencia de la historia que la histéresis de \hyperref[rep:II.1]{II.1}.
\item \textbf{\autor{Krugman} (1979)}: diferenciación + escala internas \vertema{3.B.6}: más variedades, más escala, márgenes ↓. Comercio \emph{intraindustrial} (\autor{Grubel}, 1967, CEE): ajuste más suave (empresas, no sectores).
\item \textbf{\autor{Melitz} (2003)}: \concepto{reasignación} hacia las más productivas; salen las que solo vendían dentro; productividad media ↑ por composición. ALC Canadá-EE.\,UU. (1989): $-12$\,\% empleo, $+15$\,\% productividad del trabajo (\autor{Trefler}, 2004). Coste concentrado en pymes no exportadoras.
\item \textbf{\concepto{Efectos dinámicos}}: sobre la \emph{tasa} de crecimiento, no solo el nivel \vertema{3.B.9}. Neoclásico: \emph{efecto acumulación a medio plazo} (\autor{Baldwin}, 1989), largo plazo igual. Endógeno: integrar bienes \emph{e ideas} = duplicar población, ↑ tasa (\autor{Rivera-Batiz}-\autor{Romer}, 1991; \autor{Romer}, 1990: efecto escala) \vertema{3.A.44}; \autor{Grossman}-\autor{Helpman} (1991): el de ventaja en I+D crece más. Inversión: de dentro (nueva escala) y de fuera, IED que «salta» el arancel \vertema{3.B.18}.
\item \textbf{Redistributivos}: \autor{Stolper}-\autor{Samuelson} (1941) en HOS \vertema{3.B.5}: ↑ $p$ relativo $\Rightarrow$ ↑ remuneración real del factor intensivo, ↓ la del otro. Las barreras benefician al factor escaso; integrarse con socios abundantes en él le hace perder $\Rightarrow$ presión política (puente a \hyperref[rep:II.2]{II.2}).
\end{itemize}

\sub{I.4. Formas superiores\label{rep:I.4}}
\begin{itemize}
\item \textbf{Mercado común}: \autor{MacDougall} (1960): el capital va a donde rinde más hasta igualar; triángulo de eficiencia; redistribuye (en destino ganan los trabajadores y pierde el capital local; en origen, al revés); simétrico con migración \vertema{3.B.18}. Comercio y factores \emph{sustitutivos} (\autor{Mundell}, 1957).
\item \textbf{Mercado único}: barreras no arancelarias, las que más pesan entre avanzadas \vertema{3.B.39}. \textbf{Unión monetaria}: áreas monetarias óptimas \vertema{3.B.16}. \textbf{Unión económica}: políticas fiscales y estructurales \vertema{3.B.19}.
\end{itemize}

% ==========================================================================
\bloque[5]{II. Determinantes de la integración\label{rep:II}}
\begin{ideaclave}
Se decide por razones que el bienestar no recoge: la geografía concentra la actividad cuando caen las barreras; la política elige socios y acuerdos que no son los que más bienestar dan.
\end{ideaclave}
\begin{itemize}
\item \textbf{2 caras de la moneda}: el efecto esperado determina \emph{ex ante} la integración; aquí, lo que los efectos no explican: \emph{dónde} se instala la actividad y \emph{por qué} esos socios.
\end{itemize}

\sub{II.1. Determinantes geográficos\label{rep:II.1}}
\begin{itemize}
\item Concentración (Grandes Lagos s.\,XIX, Ruhr, Silicon Valley): causas de \emph{primera naturaleza} (recursos, clima, puertos) y de \emph{segunda} (rendimientos crecientes, costes de transporte).
\item \concepto{NGE} (\autor{Fujita}, 1988; \autor{Krugman}, 1991; \autor{Venables}, 1996; FKV, 1999): pequeña asimetría $\Rightarrow$ reparto muy desigual por 2 fuerzas: \emph{aglomeración} (centrípetas: mercado amplio, productividad) y \emph{dispersión} (centrífugas: transporte, congestión, factores inmóviles). 2 sectores: tradicional (CRS) y diferenciado (IRS). Resultados muy sensibles a supuestos.
\item \textbf{Centro-periferia (\autor{Krugman}, 1991)}: como \autor{Krugman} (1979) da comercio sin ventaja comparativa, aquí concentración sin diferencias de partida. 2 regiones idénticas; agricultura CRS con trabajadores inmóviles; manufacturas \autor{Dixit}-\autor{Stiglitz} \vertema{3.A.18}, IRS, trabajadores que emigran al salario real más alto; gasto $\mu$ en manufacturas; iceberg $\tau>1$.
\item Aglomeración: \emph{mercado doméstico} (causación circular) y \emph{coste de la vida} (más variedades, $P$ ↓, salario real ↑). Dispersión: \emph{competencia local} y demanda de agricultores inmóviles.
\item \concepto{Diagrama \emph{tomahawk}} en $\phi_\tau=\tau^{1-\sigma}\in[0,1]$:
  \begin{formula}
  \[\phi_\tau=\tau^{1-\sigma},\qquad \dot\lambda_A=\omega_A(\lambda_A)-\omega_B(\lambda_A),\quad \omega_i=w_i/P_i^{\mu}\]
  \end{formula}
  $\phi_\tau$ bajo: simétrico estable; al integrarse se cruza el \emph{punto de ruptura} (simétrico inestable $\Rightarrow$ concentración); entre \emph{punto de sostenibilidad} y ruptura, 2 equilibrios: \emph{histéresis}.
\end{itemize}
\graficorepaso[0.7\linewidth]{tomahawk}{\emph{Tomahawk}: cuota de $A$ en la industria según la libertad comercial; continuas, estables}
\begin{itemize}
\item \emph{Implicaciones}: «U» de \autor{Krugman}-\autor{Venables} (1995) \vertema{3.B.40}: con costes intermedios, centro; muy bajos, la periferia (salarios bajos) atrae. La integración decide \emph{dónde} se produce $\Rightarrow$ políticas de cohesión. Completo en \vertema{3.B.6}.
\item \autor{Baldwin}-\autor{Forslid} (2000): + \autor{Romer}: comercio de \emph{bienes} $\Rightarrow$ aglomeración; de \emph{ideas} $\Rightarrow$ dispersión.
\end{itemize}

\sub{II.2. Determinantes políticos\label{rep:II.2}}
\begin{itemize}
\item Lo primero: quiénes son los socios, a menudo por motivos políticos (firma un político).
\item \textbf{5 determinantes}:
  \begin{enumerate}
  \item \emph{Poder de negociación}: el bloque negocia mejor; interesa más a los pequeños (la UE como un solo miembro) \vertema{3.B.41}.
  \item \emph{Paz y seguridad}: Francia-Alemania; lo económico fue el medio.
  \item \emph{Grupos de presión}: demanda de protección vs. oferta del gobierno (\autor{Rodrik}, 1995). \autor{Grossman}-\autor{Helpman} (1995): un ALC es más viable cuando \emph{desvía} (los productores del socio ganan cuota a costa de terceros). El riesgo de \autor{Viner} es un incentivo; sectores sensibles fuera o con transición larga.
  \item \emph{Afinidad institucional}: ASA (1961) paralizada por Sabah; ASEAN (1967), consenso y no injerencia.
  \item \emph{Efecto dominó} (\autor{Baldwin}, 1993; «bola de nieve»): los excluidos presionan para entrar.
  \end{enumerate}
\item \textbf{¿Peldaño o escollo?} (\autor{Bhagwati}, 1991). \emph{Peldaño}: dominó; entre pocos se llega más lejos y sirve de modelo en la OMC; bloques «naturales» (\autor{Frankel}-\autor{Stein}-\autor{Wei}, 1995). \emph{Escollo}: \concepto{plato de espaguetis} de reglas de origen (\autor{Bhagwati}, 1995), daña a las pymes; los ganadores de la preferencia frenan su extensión; bloques grandes elevan el arancel ($\approx$3 bloques, mínimo).
\item \textbf{Integración profunda}: con las cadenas globales de valor se protege la cadena; acuerdos de nueva generación: «regionalismo del siglo XXI» (\autor{Lawrence}, 1996; \autor{Baldwin}, 2011) \vertema{3.B.18}.
\end{itemize}

% ==========================================================================
\bloque[8]{III. Metodología empírica y evidencia\label{rep:III}}
\begin{ideaclave}
Medir es comparar con un contrafactual no observado; las herramientas convergen en modelos que dan el bienestar a partir de la gravedad, y los grandes acuerdos han creado mucho más comercio del que han desviado.
\end{ideaclave}
\begin{itemize}
\item \textbf{Idea}: \concepto{MEGC} (\emph{ex ante}, simulan) y \concepto{ecuación de gravedad} (\emph{ex post}, estima); desde 2012 convergen (\autor{Arkolakis} et al., 2012).
\end{itemize}

\sub{III.1. Modelos de equilibrio general computable\label{rep:III.1}}
\begin{itemize}
\item Walrasianos numéricos, muchos sectores, factores y países, calibrados con una MCS \vertema{3.A.21}; con y sin acuerdo.
\item CE evalúa cada acuerdo \vertema{3.B.41}: UE-Mercosur (GTAP dinámico), PIB UE +77.600 M€ en 2040, exportaciones a la región +39\,\%. España: \autor{Bajo-Rubio}-\autor{Gómez-Plana} (2005), mercado único en servicios y contratación pública: ganancias con contracción sectorial.
\item \emph{Valoración}: dependen de elasticidades y cierres; infravaloran el comercio (TLCAN, \autor{Kehoe}, 2005).
\end{itemize}

\sub{III.2. La ecuación de gravedad\label{rep:III.2}}
\begin{itemize}
\item Comercio ↑ con el tamaño, ↓ con la distancia (\autor{Tinbergen}, 1962). Nació ateórica; fundamentos: \autor{Anderson} (1979, CES por origen), \autor{Eaton}-\autor{Kortum} (2002, ricardiano).
\item \emph{Versión empírica} ($\mathrm{ACR}_{ij}=1$ si hay acuerdo):
  \funcionclave{\ln X_{ij}=\beta_0+\beta_1\ln Y_i+\beta_2\ln Y_j-\beta_3\ln d_{ij}+\beta_4\,\mathrm{ACR}_{ij}+u_{ij}}
  con acuerdo, el comercio es $e^{\beta_4}$ veces el que habría sin él.
\end{itemize}
\graficorepaso[0.6\linewidth]{gravedad}{Gravedad: comercio corregido por tamaño frente a distancia, sin y con ACR (esquema)}
\begin{itemize}
\item \emph{Versión estructural} (\autor{Anderson}-\autor{van Wincoop}, 2003):
  \begin{formula}
  \[X_{ij}=\frac{Y_iY_j}{Y_W}\left(\frac{\tau_{ij}}{\Pi_i\,P_j}\right)^{1-\sigma}\]
  \end{formula}
  \concepto{resistencias multilaterales} $\Pi_i$, $P_j$: cuenta la distancia \emph{relativa} al resto de socios; efectos fijos de exportador e importador; PPML (\autor{Santos Silva}-\autor{Tenreyro}, 2006) por flujos nulos y heterocedasticidad.
\item Con \autor{Melitz}: la distancia reduce sobre todo el margen \emph{extensivo} (\autor{Chaney}, 2008; HMR, 2008).
\end{itemize}
\begin{itemize}
\item \emph{Implicaciones}: renta $\approx$1, distancia $\approx-0{,}9$ (\autor{Head}-\autor{Mayer}, 2014, $>$2.500 coeficientes). \emph{Efecto frontera}: provincias canadienses ×22 (\autor{McCallum}, 1995); con resistencias, la frontera resta un 44\,\% (AvW). ALC: duplica el comercio en $\approx$10 años, tratando el acuerdo como endógeno (\autor{Baier}-\autor{Bergstrand}, 2007). España: \autor{Sanz} (2000), entrada en la CE: importaciones ↑ más que exportaciones.
\item \emph{Valoración}: mide bien flujos, sencilla; no da el bienestar y depende del trato de la endogeneidad.
\end{itemize}

\sub{III.3. Nuevos modelos cuantitativos\label{rep:III.3}}
\begin{itemize}
\item Gravedad = cuánto comercio; estos, cuánto \emph{bienestar}. \autor{Arkolakis}-\autor{Costinot}-\autor{Rodríguez-Clare} (2012): en \autor{Armington}, \autor{Krugman}, \autor{Eaton}-\autor{Kortum} y \autor{Melitz}, la ganancia depende solo de 2 estadísticos: \emph{cuota de gasto nacional} y \emph{elasticidad del comercio}. EE.\,UU.: 0,7-1,4\,\% de la renta real frente a autarquía. \autor{Costinot}-\autor{Rodríguez-Clare} (2014): positivas en todos, mayores con varios sectores, intermedios y competencia imperfecta. Miden el comercio \emph{existente} frente a autarquía.
\end{itemize}

\sub{III.4. Evidencia\label{rep:III.4}}
\begin{itemize}
\item \textbf{Mercado único}: \emph{ex ante}, \autor{Cecchini} (1988): PIB +4,5\,\% (3,2-5,7) a medio plazo, frente a $\approx$1\,\% de la UA de 1968; \emph{ex post}, QUEST: PIB UE +8-9\,\% (\autor{in~'t Veld}, 2019), más en pequeñas y abiertas.
\item \textbf{ALC}: comercio ×2 en una década; productividad por reasignación (\autor{Trefler}).
\item \textbf{Desintegración, \emph{Brexit}}: OBR, productividad $-4$\,\% y comercio $-15$\,\% a largo plazo; PIB 2025 entre un 6 y un 8\,\% menor (\autor{Bloom} et al., 2025): deshacerla cuesta.
\end{itemize}

% ==========================================================================
\bloque[1,5]{Conclusión}
\begin{itemize}
\item \textbf{Recapitulación}: eliminación discriminatoria; estáticos ambiguos y dominados por \autor{Cooper}-\autor{Massell}; la justifican compensaciones, RRI, escala, competencia, crecimiento; costes: factores (\autor{Stolper}-\autor{Samuelson}, \autor{MacDougall}), sectores (más si interindustrial), regiones; ciclo transmitido, desbordamientos, menos autonomía. Geografía y política; gravedad y modelos cuantitativos.
\item \textbf{Relevancia}: segundo óptimo $\Rightarrow$ cuestión empírica, no teorema; vía principal de liberalización (art.~XXIV; 390 ACR).
\item \textbf{Extensiones}: integración monetaria \vertema{3.B.16}, factores \vertema{3.B.18}, coordinación \vertema{3.B.19}; fuera de la UE \vertema{3.B.35}, mercado único \vertema{3.B.39}, política comercial UE \vertema{3.B.41}.
\item \textbf{Opinión}: la teoría estática dice que puede empeorar; la evidencia, que los grandes acuerdos crearon mucho más de lo que desviaron (escala, competencia, reasignación). Riesgo hoy: fragmentación en bloques con reglas incompatibles.
\item \textbf{Idea final}: ¿peldaño o escollo? Bloque con arancel bajo y abierto (\autor{Kemp}-\autor{Wan}, dominó) empuja; el que protege ineficientes lo frena. «Dime con quién andas\ldots»: con quién y para qué.
\end{itemize}

% ==========================================================================
\bloque{Autores y fechas}
\hito{1941}{\autor{Stolper}-\autor{Samuelson}: precios relativos y factores}
\hito{1950}{\autor{Viner}: creación y desviación de comercio}
\hito{1954}{\autor{Tinbergen}: integración negativa y positiva}
\hito{1955}{\autor{Meade}: expansión del comercio; 1956, \autor{Lipsey}-\autor{Lancaster}: segundo óptimo}
\hito{1957}{\autor{Lipsey}: creación de consumo; \autor{Mundell}: comercio y factores sustitutivos}
\hito{1960}{\autor{Lipsey}: regla; \autor{MacDougall}: capital; 1961, \autor{Balassa}: formas}
\hito{1962}{\autor{Tinbergen}: gravedad; 1964, \autor{Mundell}: país grande}
\hito{1965}{\autor{Cooper}-\autor{Massell}: rebaja no discriminatoria}
\hito{1968}{UA de la CEE completa (1/7); 1967, \autor{Grubel}: intraindustrial}
\hito{1972}{\autor{Corden}: escala; 1976, \autor{Kemp}-\autor{Wan}}
\hito{1979}{\autor{Krugman}: diferenciación; \autor{Anderson}: gravedad CES; cláusula de habilitación}
\hito{1981}{\autor{Wonnacott}-\autor{Wonnacott}; 1983, \autor{Brander}-\autor{Krugman}}
\hito{1988}{\autor{Cecchini}: mercado único; 1989, \autor{Baldwin}: acumulación}
\hito{1991}{\autor{Krugman}: centro-periferia y nº de bloques; \autor{Rivera-Batiz}-\autor{Romer}; \autor{Bhagwati}: peldaño o escollo}
\hito{1993}{\autor{Baldwin}: dominó; 1994, Entendimiento art.~XXIV}
\hito{1995}{\autor{Grossman}-\autor{Helpman}; \autor{Krugman}-\autor{Venables}: «U»; \autor{McCallum}: frontera}
\hito{1996}{\autor{Bhagwati}-\autor{Panagariya}: contra socios naturales}
\hito{2003}{\autor{Melitz}; \autor{Anderson}-\autor{van Wincoop}: resistencias multilaterales}
\hito{2006}{\autor{Santos Silva}-\autor{Tenreyro}: PPML; 2007, \autor{Baier}-\autor{Bergstrand}}
\hito{2012}{\autor{Arkolakis}-\autor{Costinot}-\autor{Rodríguez-Clare}}
\hito{2026}{UE-Mercosur, aplicación provisional (1/5)}

% ==========================================================================
\bloque{Datos clave}
\begin{itemize}
\item ACR notificados y en vigor: \dato{390, $\approx$9 de cada 10 ZLC; solo 16 UA}{OMC, nov. 2026}; sin notificar \dato{$\geq$84}{OMC}.
\item ACR en vigor: \dato{44 hasta 1990; +145 en los 90}{OMC, RTA-IS}; 2021: \dato{46 nuevos, 38 del Reino Unido}{OMC}.
\item UE-Mercosur: \dato{PIB UE +77.600 M€ en 2040; export. +39\,\%}{CE, 2025}.
\item Mercado único: \dato{+4,5\,\% PIB \emph{ex ante}}{Cecchini, 1988}; \dato{+8-9\,\% \emph{ex post}}{in 't Veld, 2019}; ALC: \dato{comercio ×2}{Baier-Bergstrand, 2007}.
\item Gravedad: \dato{distancia $\approx-0{,}9$, renta $\approx$1}{Head-Mayer, 2014}; frontera \dato{×22}{McCallum, 1995}, \dato{$-44$\,\%}{AvW, 2003}.
\item Ganancia del comercio EE.\,UU.: \dato{0,7-1,4\,\% renta real}{Arkolakis et al., 2012}.
\item \emph{Brexit}: \dato{PIB $-6$ a $-8$\,\% en 2025}{Bloom et al., 2025}; \dato{productividad $-4$\,\%, comercio $-15$\,\%}{OBR, 2026}.
\end{itemize}

% ==========================================================================
\bloque{Preguntas del tribunal}
\begin{enumerate}
\item \textbf{¿Puede la desviación mejorar el bienestar?} Con demanda fija, no. En equilibrio general (\autor{Lipsey}), sí: suprime la distorsión del consumo.
\item \textbf{¿Por qué UA si \autor{Cooper}-\autor{Massell}?} Sus supuestos: con socio con aranceles, acceso a su mercado (\autor{Wonnacott}); país grande (RRI); escala, competencia, crecimiento; motivos políticos.
\item \textbf{¿Socios naturales?} Hipótesis (\autor{Wonnacott}-\autor{Lutz}, \autor{Krugman}, \autor{Summers}) vs. \autor{Bhagwati}-\autor{Panagariya}: con socio caro, más comercio previo = más recaudación perdida; \autor{Krishna}: no predicen.
\item \textbf{¿Por qué resistencias multilaterales?} Sin ellas, sesgo (\autor{McCallum} ×22); cuenta la distancia relativa; efectos fijos y PPML.
\item \textbf{¿Peldaño o escollo?} Depende del bloque: arancel bajo y abierto (dominó) vs. espaguetis y grupos de presión.
\item \textbf{¿Por qué es viable un ALC que desvía?} \autor{Grossman}-\autor{Helpman}: los productores del socio ganan cuota a costa de terceros y lo apoyan.
\end{enumerate}

\begin{trampa}
\begin{itemize}
\item Desviación del \emph{tráfico} (ZLC, reglas de origen) $\neq$ desviación de \emph{comercio} (\autor{Viner}).
\item \autor{Viner}: demanda \emph{perfectamente inelástica}, costes constantes; $\Delta W=(p_R-p_B)\bar x<0$ en la desviación.
\item \autor{Lipsey}: $\Delta W=b+d-e$; $c$ es transferencia (de recaudación a consumidores), no ganancia neta.
\item \autor{Cooper}-\autor{Massell}: $b+d+f+g>0$ \emph{siempre}; la recaudación puede subir o bajar ($f+g-c$).
\item \autor{Corden}: reducción de costes (+) en $B$; supresión de comercio ($-$) cuando nace una empresa; caso a) positivo, b) ambiguo, c) negativo.
\item \emph{Tomahawk}: ruptura y sostenibilidad; histéresis entre ambos; $\phi_\tau=\tau^{1-\sigma}$ va de 0 (imposible) a 1 (libre).
\item \autor{Baldwin}-\autor{Forslid}: ideas $\Rightarrow$ \emph{dispersión}; bienes $\Rightarrow$ aglomeración.
\item Art.~XXIV: ZLC y UA; preferencias parciales, solo entre PED (habilitación, 1979); servicios, art.~V AGCS.
\item Gravedad: $\beta_4$ $\Rightarrow$ comercio $e^{\beta_4}$ veces; distancia reduce el margen \emph{extensivo}; \autor{Arkolakis} et al. miden el comercio existente frente a autarquía.
\end{itemize}
\end{trampa}

\end{columnas}
\end{document}
